% Source: https://artofproblemsolving.com/community/c3h2456029_mathcounts_2009_state_sprint_round_problem_29
% Stable identity: aops-topic-2456029-q01
Take a list with sum $22$ whose product is as large as possible (there are finitely many lists, so one exists). We improve any list that has one of the following features, without changing the sum:
\begin{itemize}
\item a part $k\ge5$: replace it by $2$ and $k-2$; the product is multiplied by $\dfrac{2(k-2)}{k}>1$, since $2k-4>k$;
\item a part $1$ (with at least one other part $c$): replace $1$ and $c$ by $c+1$, which is larger than $1\cdot c$;
\item three parts equal to $2$: replace them by $3,3$, since $9>8$.
\end{itemize}
So a list with the largest product has no part $1$ and no part $5$ or more, and at most two parts equal to $2$. Replacing each part $4$ by $2,2$ keeps the sum and the product, so it gives another list with the largest product, which by the same reasoning has at most two $2$'s. Hence some list with the largest product consists only of $3$'s and at most two $2$'s. Since $22=3\cdot6+2\cdot2$ and $22$ leaves remainder $1$ when divided by $3$, the only such list is six $3$'s and two $2$'s, with product
\[3^6\cdot2^2=729\cdot4=\boxed{2916}.\]
